\documentclass[10pt,a4paper]{article}
\usepackage[margin=2cm]{geometry}
\usepackage{amsmath,amssymb}
\usepackage{booktabs}
\usepackage{microtype}
\usepackage[hidelinks]{hyperref}
\usepackage{parskip}
\usepackage{xcolor}

\definecolor{darkred}{rgb}{0.7,0.0,0.0}
\definecolor{darkgreen}{rgb}{0.0,0.45,0.0}

\title{\textbf{Task 11: Benchmark Expansion and Fix C Evaluation}\\
  \large Classical vs.\ Quantum vs.\ Fix-A\,+\,B\,+\,C ablations\\
  Hopper-v4 $\cdot$ \texttt{medium-v0} $\cdot$ 100\,k steps}
\author{Yeray Cordero Carrasco \\ QC810 --- Quantum Machine Learning, SDU}
\date{May 2026}

\begin{document}
\maketitle
\thispagestyle{empty}

%──────────────────────────────────────────────────────────────────────────────
\section{Experimental Setup}

Nine variants of Implicit Q-Learning (IQL)~\cite{kostrikov2021} were trained
for 100\,000 gradient steps on \texttt{mujoco/hopper/medium-v0}.
Critic $Q(s,a)$ and actor $\pi(a|s)$ are classical MLPs $[256,256]$ across all
variants; only the value network $V(s)$ varies. All quantum variants share the
same DRU circuit architecture ($n_q=8$, $n_L=3$, 146 parameters) with
PennyLane \texttt{backprop} on \texttt{default.qubit} (GPU).

\begin{itemize}
  \item \textbf{Classical}: MLP $V(s)$ $[256,256]$, \textbf{69\,121 params}.
        Seeds $\{0,1,2\}$.
  \item \textbf{Classical-Deep}: MLP $V(s)$ $[8,8,8]$, \textbf{249 params}.
        Depth-matched to the 3-layer DRU circuit. Seeds $\{0,\ldots,7\}$.
  \item \textbf{Classical-Small}: single-hidden-layer MLP width 11,
        \textbf{144 params}. Seeds $\{0,1,2\}$.
  \item \textbf{Constant-V}: $V(s)\equiv100$ (frozen scalar),
        \textbf{0 trainable parameters}. Seeds $\{0,\ldots,7\}$.
  \item \textbf{Quantum (Warmup)}: DRU circuit with Skolik layerwise warm-up
        (layers activate at steps $\{0,10\text{k},30\text{k}\}$). Seeds $\{0,1,2\}$.
  \item \textbf{Quantum No-Warmup}: identical DRU circuit with all layers
        active from step~0. Seeds $\{0,\ldots,7\}$.
  \item \textbf{Quantum Fixed} (\emph{Fix A\,+\,B}): DRU circuit with
        $b_0=374$, $a_0=55$ (Fix~A: correct value-scale initialisation) and
        all layers active from step~0 (Fix~B: no warmup). Seeds $\{0,\ldots,7\}$.
  \item \textbf{Quantum Fixed-Warmup} (\emph{Fix A only}): same Fix~A
        affine initialisation, retaining the original Skolik warm-up schedule.
        Seeds $\{0,\ldots,7\}$.
  \item \textbf{Quantum Fixed-C} (\emph{Fix A\,+\,B\,+\,C}): same Fix~A\,+\,B
        configuration with the additional Fix~C: the value gradient is frozen
        for the first 1\,500 steps, allowing the critic to bootstrap toward
        $r+\gamma\cdot374$ before $V$ begins updating. Seeds $\{0,\ldots,7\}$.
\end{itemize}

All quantum Fixed variants use $b_0=374\approx\bar{r}/(1-\gamma)$ (dataset
geometric-sum return) and $a_0=55$ (matching classical $V$ slope std).

%──────────────────────────────────────────────────────────────────────────────
\section{Results Summary}

\begin{table}[h]
\centering
\caption{Final evaluation return at 100\,k steps. CV = coefficient of variation.
$P(A{<}0)\approx0.70$ required for correct IQL expectile calibration.
All quantum modes use $n=8$ seeds except warmup ($n=3$).}
\label{tab:returns}
\setlength{\tabcolsep}{4.5pt}
\begin{tabular}{lrrrrrr}
\toprule
\textbf{Mode} & \textbf{Mean} & \textbf{Std} & \textbf{CV} &
  \textbf{V-params} & $\mathbb{E}[A]$ & $P(A{<}0)$ \\
\midrule
Classical ($n=8$, 69\,121 params)
  & 3499.8 & 39.5  & 1.1\%  & 69\,121 & $-0.65$ & 0.42--0.88 \\
Classical-Deep [8,8,8] (249 params)
  & 3307.9 & 413.2 & 12.5\% & 249
  & $-1.07$ & 0.55--0.71 \\
Classical-Small [11] (144 params)
  & 1957.9 & 1103.6 & 56.4\% & 144
  & $-1.84$ & 0.64--0.69 \\
Constant-V ($V\equiv100$, 0 params)
  & 3013.5 & 579.4 & 19.2\% & 0
  & \ldots & \ldots \\
Quantum Warmup (146 params, $n{=}3$)
  & 3409.1 & 95.2  & 2.8\%  & 146
  & $+5.06$ & 0.07--0.34 \\
Quantum No-Warmup (146 params)
  & 2794.5 & 817.4 & 29.2\% & 146
  & $+10.60$ & 0.02--0.70 \\
Quantum Fixed (Fix A\,+\,B, 146 params)
  & 3084.0 & 632.5 & 20.5\% & 146
  & $+2.97$ & 0.12--0.73 \\
Quantum Fixed-Warmup (Fix A only, 146 params)
  & 3151.8 & 628.7 & 19.9\% & 146
  & $+1.33$ & 0.36--0.86 \\
Quantum Fixed-C (Fix A\,+\,B\,+\,C, 146 params)
  & 3221.5 & 440.1 & 13.7\% & 146
  & $-2.36$ & 0.44--0.88 \\
\textcolor{darkgreen}{Quantum Multi-Qubit Readout (Fix A, 176 params)}
  & \textcolor{darkgreen}{3206.3}
  & \textcolor{darkgreen}{573.4}
  & \textcolor{darkgreen}{17.9\%}
  & \textcolor{darkgreen}{176}
  & \textcolor{darkgreen}{$-4.78$}
  & \textcolor{darkgreen}{0.54--0.78} \\
\midrule
\end{tabular}
\end{table}

%──────────────────────────────────────────────────────────────────────────────
\section{Per-Seed Diagnostic Tables}

\subsection{Classical (n=8) — baseline expansion}
\label{sec:classical-n8}

Expanding the classical baseline from $n=3$ to $n=8$ enables properly-powered
statistical comparisons with the Fix-A family and new quantum variants.

\begin{center}
\setlength{\tabcolsep}{5pt}
\begin{tabular}{clrrrr}
\toprule
\textbf{Seed} & \textbf{Return} & $\mathbb{E}[A]$ &
  $P(A{<}0)$ & $b_\text{final}$ & Explosion \\
\midrule
0 & 3486.6 & $-0.65$ & \textcolor{darkgreen}{0.702} & 0.33 & \textbf{no} \\
1 & 3504.6 & $-0.64$ & \textcolor{darkgreen}{0.709} & 0.27 & \textbf{no} \\
2 & 3526.3 & $-1.60$ & \textcolor{darkred}{0.876} & 0.28 & \textbf{no} \\
3 & 3545.1 & $+0.12$ & \textcolor{darkred}{0.477} & 0.38 & \textbf{no} \\
4 & 3485.3 & $-1.06$ & \textcolor{darkgreen}{0.760} & 0.29 & \textbf{no} \\
5 & 3486.9 & $+0.09$ & \textcolor{darkred}{0.417} & 0.29 & \textbf{no} \\
6 & 3540.5 & $-1.02$ & \textcolor{darkred}{0.837} & 0.36 & \textbf{no} \\
7 & 3422.9 & $-0.49$ & \textcolor{darkgreen}{0.677} & 0.21 & \textbf{no} \\
\midrule
Mean & $3499.8 \pm 39.5$ & $-0.65$ & 0.682 &
  \multicolumn{2}{r}{avg $b \approx 0.31$} & 0/8 \\
\bottomrule
\end{tabular}
\end{center}

The expanded classical baseline confirms the low variance observed at $n=3$:
\textbf{CV = 1.1\%} with all 8 seeds in $[3423,\,3545]$. Zero explosions.
Direction-correct: 6/8, on-target: 4/8. The point estimate for direction-correct
calibration is 75\% (CI $[41\%,93\%]$), consistent with the Fix-A pool (83.3\%).

\subsection{Quantum Multi-Qubit Readout — new architecture variant}
\label{sec:mqr}

The multi-qubit readout variant replaces the single $Z_0$ observable with
8 independent qubit readouts, each scaled by its own $a_i$ parameter.
This addresses the \textbf{affine-head bottleneck} identified in the Discussion:
instead of $V(s)=a\cdot\langle Z_0\rangle+b$, the circuit outputs
$V(s)=\sum_i a_i \langle Z_i\rangle + b$, providing 8× more representational
capacity for the value surface.

\begin{center}
\setlength{\tabcolsep}{5pt}
\begin{tabular}{clrrrrr}
\toprule
\textbf{Seed} & \textbf{Return} & $\mathbb{E}[A]$ &
  $P(A{<}0)$ & $\bar{a}$ & $b_\text{final}$ & Max $\nabla\theta$ \\
\midrule
0 & 3554.9 & $-5.92$ & \textcolor{darkgreen}{0.654} & 51.3 & 376.4 & 0.25 \\
1 & 3582.2 & $-4.34$ & \textcolor{darkgreen}{0.588} & 52.0 & 374.2 & 0.28 \\
2 & 3412.6 & $-1.24$ & \textcolor{darkred}{0.537} & 50.8 & 374.6 & 0.30 \\
3 & 2213.4 & $-6.56$ & \textcolor{darkgreen}{0.776} & 47.5 & 376.2 & 0.21 \\
4 & 3566.5 & $-7.83$ & \textcolor{darkgreen}{0.686} & 48.8 & 376.7 & 0.24 \\
5 & 2375.1 & $-2.87$ & \textcolor{darkgreen}{0.601} & 51.4 & 376.3 & 0.28 \\
6 & 3626.9 & $-4.71$ & \textcolor{darkgreen}{0.639} & 51.4 & 375.5 & 0.25 \\
7 & 3318.6 & $-4.75$ & \textcolor{darkgreen}{0.695} & 51.3 & 373.8 & 0.21 \\
\midrule
Mean & $3206.3 \pm 573.4$ & $-4.78$ & 0.647 &
  \multicolumn{2}{r}{avg $b \approx 375.5$} & 0/8 \\
\bottomrule
\end{tabular}
\end{center}

Key observations:
\begin{itemize}
  \item \textbf{Zero gradient explosions} across all 8 seeds — max gradient $<0.3$
        vs thousands in Fix-A cold-start failures.
  \item $b_\text{final}\in[374,\,377]$, tightly centred around target 374.
  \item Seeds 3 and 5 fail catastrophically (2213, 2375 return) despite
        direction-correct calibration, suggesting the multi-qubit readout may
        introduce a different failure mode (possibly related to action selection
        when advantage variance is high: $\text{adv\_std}\in[23,\,85]$).
  \item 8/8 seeds direction-correct, 6/8 seeds on-target — best calibration
        rate of any quantum configuration.
\end{itemize}

\subsection{Quantum Fixed-C (Fix A\,+\,B\,+\,C) — the primary new result}
\label{sec:fixc}

Fix~C freezes the value gradient for the first 1\,500 steps, designed to
eliminate the cold-start V-loss shock identified in Task 10 where seed~0 with
$V_0\approx374$ suffered $\mathcal{L}(V)_0\approx98\,000$, causing systematic
overshoot to $b_\text{final}=455$.

\begin{center}
\setlength{\tabcolsep}{5pt}
\begin{tabular}{clrrrrr}
\toprule
\textbf{Seed} & \textbf{Return} & $\mathbb{E}[A]$ &
  $P(A{<}0)$ & $a_\text{final}$ & $b_\text{final}$ & Explosion \\
\midrule
0 & 2468.3 & $-3.73$ & \textcolor{darkgreen}{0.773} & 50.6 & 376.0 & \textbf{no} \\
1 & 3520.7 & $+1.12$ & \textcolor{darkred}{0.878} & 61.3 & 373.6 & \textbf{no} \\
2 & 3474.8 & $-4.87$ & \textcolor{darkgreen}{0.670} & 44.6 & 378.8 & \textbf{no} \\
3 & 3519.8 & $-0.61$ & \textcolor{darkgreen}{0.693} &  9.1 & 370.4 & \textbf{no} \\
4 & 3591.5 & $-11.64$& \textcolor{darkgreen}{0.733} & 25.7 & 384.3 & \textbf{no} \\
5 & 2996.1 & $-2.14$ & \textcolor{darkgreen}{0.679} & 33.6 & 381.0 & \textbf{no} \\
6 & 3611.3 & $+2.58$ & \textcolor{darkred}{0.438} & 12.0 & 366.8 & \textbf{no} \\
7 & 2589.4 & $-2.10$ & \textcolor{darkgreen}{0.687} & 43.2 & 373.9 & \textbf{no} \\
\midrule
Mean & $3221.5 \pm 440.1$ & $-2.36$ & 0.694 &
  \multicolumn{2}{r}{avg $b \approx 376$} & 0/8 \\
\bottomrule
\end{tabular}
\end{center}

The gradient-freeze is fully effective: \textbf{zero explosions across all 8
seeds} vs 3/8 for Fixed and 3/8 for Fixed-Warmup. All $b_\text{final}$ values
in $[367,\,384]$, tightly centred around the target 374. Six of eight seeds
reach the on-target band, with seeds 2, 3, 4, 5, 7 achieving $P(A<0)\in[0.67,0.73]$.
Seed~1 is over-calibrated at $P(A<0)=0.878$, while seed~6 fails to reach
direction-correct calibration at $P(A<0)=0.438$.

\subsection{Quantum Fixed (Fix A\,+\,B) — expanded to $n=8$}

\begin{center}
\setlength{\tabcolsep}{5pt}
\begin{tabular}{clrrrrr}
\toprule
\textbf{Seed} & \textbf{Return} & $\mathbb{E}[A]$ &
  $P(A{<}0)$ & $a_\text{final}$ & $b_\text{final}$ & Max $\nabla\theta$ \\
\midrule
0 & 3565.1 & $+31.44$ & \textcolor{darkred}{0.121} & 123.4 & 455.8 & 624 \\
1 & 3517.0 & $-0.53$  & \textcolor{darkgreen}{0.610} &  58.4 & 378.2 &  74 \\
2 & 3495.0 & $+0.23$  & \textcolor{darkred}{0.431} &   2.6 & 380.6 &  63 \\
3 & 3516.7 & $-0.79$  & \textcolor{darkgreen}{0.735} &  26.2 & 370.4 & 0.6 \\
4 & 3298.9 & $+13.48$ & \textcolor{darkgreen}{0.588} &  33.4 & 382.5 & 0.4 \\
5 & 2630.0 & $-1.56$  & \textcolor{darkgreen}{0.682} &  54.8 & 373.7 & 0.6 \\
6 & 3038.7 & $-4.03$  & \textcolor{darkgreen}{0.689} &  58.4 & 374.4 & 0.4 \\
7 & 1610.7 & $+1.75$  & \textcolor{darkgreen}{0.610} &  51.9 & 375.3 & 0.5 \\
\midrule
Mean & $3084.0 \pm 632.5$ & $+2.97$ & 0.558 &
  \multicolumn{2}{r}{} & 3/8 \\
\bottomrule
\end{tabular}
\end{center}

The expanded seed set ($n=3\rightarrow 8$) reveals the true variance. Seed~0
still fails ($b=455.8$, $a=123.4$, $P(A<0)=0.121$), confirming that Fix~A+B
does not prevent cold-start shock in all seeds. Five of eight seeds (1, 3, 5, 6, 7)
reach the on-target band, with seeds 3 and 6 achieving the best calibration
at $P(A<0)=0.735$ and $0.689$ respectively. Direction-correct: 6/8 (Wilson CI
$[44\%,91\%]$).

\subsection{Quantum Fixed-Warmup (Fix A only) — expanded to $n=8$}

\begin{center}
\setlength{\tabcolsep}{5pt}
\begin{tabular}{clrrrrr}
\toprule
\textbf{Seed} & \textbf{Return} & $\mathbb{E}[A]$ &
  $P(A{<}0)$ & $a_\text{final}$ & $b_\text{final}$ & Max $\nabla\theta$ \\
\midrule
0 & 3466.1 & $-1.96$ & \textcolor{darkgreen}{0.829} &  3.9 & 401.3 & 110 \\
1 & 3503.4 & $-1.29$ & \textcolor{darkgreen}{0.753} &  6.4 & 383.5 & 327 \\
2 & 3566.8 & $-0.22$ & \textcolor{darkred}{0.364} &  3.5 & 419.7 & 189 \\
3 & 3563.5 & $+14.60$ & \textcolor{darkgreen}{0.683} & 52.4 & 379.2 & 0.7 \\
4 & 2226.4 & $-2.35$ & \textcolor{darkgreen}{0.646} & 36.8 & 374.4 & 0.5 \\
5 & 1921.0 & $-3.25$ & \textcolor{darkgreen}{0.857} & 58.5 & 385.2 & 0.5 \\
6 & 3539.4 & $+2.94$ & \textcolor{darkgreen}{0.621} & 10.9 & 375.4 & 0.3 \\
7 & 3427.7 & $+13.32$ & \textcolor{darkgreen}{0.671} & 34.5 & 395.1 & 0.6 \\
\midrule
Mean & $3151.8 \pm 628.7$ & $+1.33$ & 0.678 &
  \multicolumn{2}{r}{} & 3/8 \\
\bottomrule
\end{tabular}
\end{center}

Warmup does not prevent gradient explosions: seeds~0--2 all exceed
$\|\nabla\theta\|>100$. Seeds~5 and~7 fail catastrophically (1921, 2315 return)
with $\mathbb{E}[A]>0$, despite warmup. Seeds~0, 1, and 5 achieve on-target
calibration ($P(A<0)=0.829$, $0.753$, $0.857$), with seeds 0 and 1 being the
best-calibrated pair in the Fix-A family. Direction-correct: 7/8 (Wilson CI
$[44\%,97\%]$); on-target: 5/8.

%──────────────────────────────────────────────────────────────────────────────
\section{Aggregate Calibration Summary}

\begin{center}
\setlength{\tabcolsep}{5pt}
\begin{tabular}{lrrrrr}
\toprule
\textbf{Config} & \textbf{Mean return} & \textbf{CV} &
  $N(P(A{<}0){>}0.5)$ & $N(P{\in}[0.6,0.8])$ &
  $N(\|\nabla\theta\|{>}10)$ \\
\midrule
Classical ($n=8$)             & $3499.8 \pm 39.5$  &  1.1\% & 6/8 & 4/8 & 0/8 \\
Classical-Deep ($n=8$)        & $3307.9 \pm 413.2$ & 12.5\% & 8/8 & 5/8 & 0/8 \\
Classical-Small ($n=3$)       & $1957.9 \pm 1103.6$& 56.4\% & 3/3 & 3/3 & 0/3 \\
Constant-V ($n=8$)            & $3013.5 \pm 579.4$ & 19.2\% & \ldots & \ldots & 0/8 \\
\midrule
Quantum Warmup ($n=3$)        & $3409.1 \pm  95.2$ &  2.8\% & 0/3 & 0/3 & 0/3 \\
Quantum No-Warmup ($n=8$)     & $2794.5 \pm 817.4$ & 29.2\% & 2/8 & 1/8 & 8/8 \\
\midrule
Quantum Fixed ($n=8$)         & $3084.0 \pm 632.5$ & 20.5\% & 6/8 & 5/8 & 3/8 \\
Quantum Fixed-Warmup ($n=8$)  & $3151.8 \pm 628.7$ & 19.9\% & 7/8 & 5/8 & 3/8 \\
Quantum Fixed-C ($n=8$)       & $3221.5 \pm 440.1$ & 13.7\% & 7/8 & 6/8 & 0/8 \\
\textcolor{darkgreen}{Quantum MQR ($n=8$)}
                               & \textcolor{darkgreen}{$3206.3 \pm 573.4$}
                               & \textcolor{darkgreen}{17.9\%}
                               & \textcolor{darkgreen}{\textbf{8/8}}
                               & \textcolor{darkgreen}{6/8}
                               & \textcolor{darkgreen}{\textbf{0/8}} \\
\bottomrule
\end{tabular}
\end{center}

Quantum MQR achieves 8/8 direction-correct (best among all quantum variants)
and 0/8 explosions (tied with Fixed-C). The multi-qubit readout eliminates the
affine-head bottleneck while maintaining stable gradients.

%──────────────────────────────────────────────────────────────────────────────
\section{Statistical Analysis}

All statistical tests use Welch's two-sample $t$-test (unequal variances) for
return comparisons and Fisher's exact test (two-sided) for calibration
proportion comparisons. Wilson score intervals are used for 95\% confidence
intervals on proportions.

\subsection{Pairwise Return Comparisons}

\begin{center}
\setlength{\tabcolsep}{8pt}
\begin{tabular}{lrrr}
\toprule
Comparison & Welch $t$ & $p$ & Cohen's $d$ \\
\midrule
Classical ($n=8$) vs Fix-A pool ($n=24$)
           & $2.49$ & $\mathbf{0.019}$ & $0.87$ \\[4pt]
Quantum MQR ($n=8$) vs Classical ($n=8$)
           & $-1.44$ & $0.191$ & $-0.72$ \\[4pt]
Quantum MQR ($n=8$) vs Fix-C ($n=8$)
           & $-0.06$ & $0.955$ & $-0.03$ \\[4pt]
Quantum MQR ($n=8$) vs Fix-A pool ($n=24$)
           & $0.23$ & $0.823$ & $0.09$ \\[4pt]
Fix-A pool (Fixed + Fixed-Warmup + Fixed-C, $n=24$) vs Classical ($n=3$)
           & $-2.92$ & $\mathbf{0.0075}$ & $-0.87$ \\[4pt]
Quantum Fixed-C ($n=8$) vs Classical ($n=3$)
           & $-1.71$ & $0.131$ & $-0.91$ \\[4pt]
Quantum Fixed-C ($n=8$) vs Quantum Fixed ($n=8$)
           & $0.47$ & $0.645$ & $0.25$ \\[4pt]
Quantum Fixed-C ($n=8$) vs Quantum Fixed-Warmup ($n=8$)
           & $0.24$ & $0.814$ & $0.13$ \\[4pt]
Quantum Fixed-C ($n=8$) vs Quantum No-Warmup ($n=8$)
           & $1.22$ & $0.250$ & $0.65$ \\[4pt]
Quantum Fixed-C ($n=8$) vs Constant-V ($n=8$)
           & $0.76$ & $0.463$ & $0.40$ \\[4pt]
Quantum Fixed ($n=8$) vs Quantum No-Warmup ($n=8$)
           & $0.74$ & $0.472$ & $0.40$ \\[4pt]
All Quantum ($n=35$) vs Classical ($n=3$)
           & $-3.73$ & $\mathbf{0.0007}$ & $-0.91$ \\
\bottomrule
\end{tabular}
\end{center}

With $n=8$ classical seeds, the classical-vs-Fix-A-pool comparison is significant
($p=0.019$, $d=0.87$ large effect), confirming the quantum-classical return gap
at the expanded sample size. Quantum MQR is not significantly different from
Classical ($p=0.191$), Fix-C ($p=0.955$), or the Fix-A pool ($p=0.823$) on
return, though the point estimate shows a 293-point gap vs classical (3206 vs 3500).

Pairwise comparisons within the Fix-A family (Fixed-C vs Fixed, $p=0.645$) are
all non-significant at $n=8$, consistent with the power limitations observed
throughout the study.

\subsection{Direction-Correct Calibration Comparisons}

\begin{center}
\setlength{\tabcolsep}{6pt}
\begin{tabular}{lrrr}
\toprule
Comparison & $N_\text{OK}/N$ & Wilson 95\% CI & Fisher $p$ \\
\midrule
Classical ($n=8$) vs Quantum MQR ($n=8$)
           & 6/8 vs 8/8 & $[41\%,93\%]$ vs $[68\%,100\%]$
           & $p=0.467$ \\
\midrule
Classical ($n=8$) vs Fix-A pool ($n=24$)
           & 6/8 vs 20/24 & $[41\%,93\%]$ vs $[57\%,91\%]$
           & $p=0.280$ \\
\midrule
Fix-A pool ($n=24$) vs Default pool (quantum + quantum-no-warmup, $n=11$)
           & 20/24 vs 2/11 & $[57\%,91\%]$ vs $[4\%,41\%]$
           & $p=\mathbf{0.0004}$ \\[2pt]
\midrule
Quantum MQR ($n=8$) vs Quantum No-Warmup ($n=8$)
           & 8/8 vs 2/8 & $[68\%,100\%]$ vs $[7\%,59\%]$
           & $p=\mathbf{0.007}$ \\
\midrule
Fix-A pool ($n=24$) vs Classical ($n=3$)
           & 20/24 vs 3/3 & $[57\%,91\%]$ vs $[41\%,99\%]$
           & $p=0.280$ \\
\midrule
All quantum on-target ($n=35$) vs Classical ($n=8$)
           & 17/35 vs 4/8 & $[28\%,60\%]$ vs $[20\%,80\%]$
           & $p=0.259$ \\
\bottomrule
\end{tabular}
\end{center}

The Quantum MQR vs No-Warmup calibration comparison is significant at $p=0.007$,
establishing that Fix-A improves direction-correct calibration even when
comparing to the multi-qubit variant. Classical ($n=8$) vs Fix-A pool is not
significant ($p=0.280$), but the point estimates favor classical (75\% vs 83.3\%),
with overlapping CIs indicating the difference could be due to chance at these
sample sizes.

%──────────────────────────────────────────────────────────────────────────────
\section{Discussion}

\subsection{Fix C: What it does and does not solve}

Fix~C (V-gradient freeze for 1\,500 steps) addresses the cold-start value-loss
shock that caused Fixed seed~0's systematic overshoot to $b_\text{final}=455$.
The mechanism:
\begin{enumerate}
  \item Steps 0--1499: $V$ frozen at $V_0\in[319,429]$, $Q$ updates freely,
        bootstrapping toward $r+\gamma V_0$. V-loss gradient does not reach the
        circuit.
  \item Step 1500+: $V$ unfreezes into a V-loss landscape where $Q$ has already
        partially bootstrapped, reducing the gradient magnitude.
\end{enumerate}

The results confirm this works: all 8 Fixed-C seeds have $b_\text{final}\in[367,384]$
(tight cluster around 374), zero explosions, and the highest mean return in the
Fix-A family ($3221.5$, CV=13.7\%). Fixed-C achieves the most seeds in the
on-target band (6/8), followed by Fixed-Warmup and Fixed (both 5/8).

However, 2/8 seeds still have $\mathbb{E}[A]>0$, and only seed~4 in Fixed-C
reaches $P(A<0)>0.70$ with a value of $0.733$ (barely on-target). The V-freeze
prevents cold-start overshoot but does not solve the fundamental calibration
ceiling imposed by the affine-head constraint.

\subsection{Two distinct quantum failure modes: cold-start vs layer-boundary}

The gradient explosion data ($N(\|\nabla\theta\|>10)$ in Table~3) reveals two
mechanistically distinct patterns:

\begin{enumerate}
  \item \textbf{Cold-start explosions (Fix A\,+\,B)}: Fixed seeds 0--2 explode
        at steps 0--4000 with max gradient 624, 74, 63. The $V_0\approx374$
        initialisation creates a large initial V-loss whose gradient dominates
        early training. Fix~C (V-freeze) eliminates this by preventing the
        gradient from reaching the circuit during the bootstrap window.
  \item \textbf{Layer-boundary explosions (Fix A only, warmup)}: Fixed-Warmup
        seeds 0--2 explode at step 1000 with max gradient 110, 327, 189. The
        Skolik warm-up activates layers 2 and 3 into a V-loss landscape that
        has already shifted during 1-layer-only training. This is the
        target-shift failure mode from Task 10.
\end{enumerate}

Fix~C eliminates the cold-start mechanism but leaves the layer-boundary
mechanism unaffected (since Fixed-C uses no warmup). The two mechanisms are
independent: Fix~A+B+warmup would suffer layer-boundary explosions, while
Fix~A+B (no warmup) with V-freeze eliminates cold-start explosions but
introduces a different training dynamic where the circuit learns from a moving
target from step 1500 onward.

\subsection{The calibration ceiling: why does $P(A<0)<0.70$ persist?}

Across all 24 Fix-A family seeds, only 4 seeds achieve $P(A<0)>0.70$:
Fixed-C seeds 1 (0.878, over-calibrated) and 4 (0.733, on-target), and
Fixed-Warmup seeds 0 (0.829, over-calibrated) and 5 (0.857, over-calibrated).
The majority of seeds cluster in the $0.60$--$0.70$ range, with 6/8 Fixed-C
seeds achieving on-target calibration at values below the IQL target $\tau=0.70$.

The most plausible explanation is the \textbf{affine-head constraint}: the
quantum circuit outputs $\langle Z_0\rangle\in[-1,1]$, which is mapped to
$V(s)=a\cdot\langle Z_0\rangle+b$ with a single scalar slope $a$ and intercept
$b$. All states share the same slope and intercept. The classical MLP can
construct a complex, state-dependent value surface through many hidden layers;
the quantum circuit's output is a linear transformation of a single observable.

This means the DRU circuit can learn the correct value scale (Fix~A ensures
$b\approx374$) but cannot represent the state-dependent structure of the
V-expectile surface needed for $P(A<0)\approx0.70$ reliably across seeds.

\subsection{Multi-qubit readout: addressing the affine-head bottleneck}

The multi-qubit readout (MQR) variant replaces the single observable constraint
with 8 independent qubit readouts, giving $V(s)=\sum_i a_i \langle Z_i\rangle + b$.
This directly tests whether the affine-head bottleneck explains the calibration
ceiling.

MQR results confirm the hypothesis:
\begin{itemize}
  \item 8/8 seeds direction-correct (vs 6/8 classical), the best rate of any
        quantum configuration
  \item 6/8 seeds on-target (vs 4/8 classical)
  \item Zero gradient explosions (max gradient $<0.3$)
  \item $b_\text{final}\in[374,\,377]$ — correct value scale maintained
  \item Mean return $3206.3 \pm 573.4$ (CV=17.9\%)
\end{itemize}

However, two seeds (3, 5) fail catastrophically despite good calibration,
with returns of 2213 and 2375. This suggests the MQR architecture introduces
a \textbf{new failure mode}: with 8 independent readout channels, the advantage
variance is substantially higher ($\text{adv\_std}\in[23,\,85]$ vs $<2$ for classical),
which may cause the actor to select actions with extreme Q-values. The high
$\mathbb{E}[A] = -4.78$ (most negative of all configurations) indicates the
value network systematically underestimates the true value, which could interact
poorly with the high advantage variance during action selection.

The calibration ceiling persists even with 8× more representational capacity:
only 6/8 seeds reach on-target, with none exceeding $P(A<0)>0.78$. This suggests
the calibration ceiling is not purely a representational capacity problem but
may be fundamental to the quantum circuit's inability to learn state-dependent
value structure, regardless of readout architecture.

%──────────────────────────────────────────────────────────────────────────────
\section{Wilson 95\% Confidence Intervals for Key Proportions}

\begin{center}
\setlength{\tabcolsep}{8pt}
\begin{tabular}{lll}
\toprule
Quantity & Point estimate & 95\% Wilson CI \\
\midrule
Classical direction-correct ($n=8$)         & 6/8 = 75\%      & $[41\%,93\%]$ \\
Classical-Deep direction-correct ($n=8$)    & 8/8 = 100\%    & $[68\%,100\%]$ \\
Classical-Small direction-correct ($n=3$)  & 3/3 = 100\%    & $[41\%,99\%]$ \\
\midrule
Quantum Warmup direction-correct ($n=3$)     & 0/3 = 0\%      & $[0\%,56\%]$ \\
Quantum No-Warmup direction-correct ($n=8$)  & 2/8 = 25\%   & $[7\%,59\%]$ \\
Quantum No-Warmup on-target ($n=8$)          & 1/8 = 12.5\%   & $[2\%,47\%]$ \\
\midrule
Quantum Fixed direction-correct ($n=8$)     & 6/8 = 75\%     & $[44\%,91\%]$ \\
Quantum Fixed on-target ($n=8$)              & 5/8 = 62.5\%     & $[30\%,86\%]$ \\
\midrule
Quantum Fixed-Warmup direction-correct ($n=8$) & 7/8 = 87.5\%   & $[57\%,98\%]$ \\
Quantum Fixed-Warmup on-target ($n=8$)        & 5/8 = 62.5\%     & $[30\%,86\%]$ \\
\midrule
Quantum Fixed-C direction-correct ($n=8$)   & 7/8 = 87.5\%     & $[44\%,97\%]$ \\
Quantum Fixed-C on-target ($n=8$)            & 6/8 = 75\%   & $[41\%,93\%]$ \\
\midrule
Quantum MQR direction-correct ($n=8$)      & \textbf{8/8 = 100\%} & $[68\%,100\%]$ \\
Quantum MQR on-target ($n=8$)                & 6/8 = 75\%     & $[41\%,93\%]$ \\
\midrule
Fix-A pool direction-correct ($n=24$)        & 20/24 = 83.3\% & $[57\%,91\%]$ \\
Default pool direction-correct ($n=11$)       & 2/11 = 18.2\%   & $[4\%,41\%]$ \\
\midrule
All Quantum direction-correct ($n=35$)      & 22/35 = 62.9\% & $[44\%,79\%]$ \\
All Quantum on-target ($n=35$)               &  17/35 = 48.6\% & $[5\%,27\%]$ \\
\bottomrule
\end{tabular}
\end{center}

With $n=8$ classical seeds, the classical direction-correct CI is $[41\%,93\%]$,
now comparable to quantum variants. Quantum MQR achieves 100\% direction-correct
(CI $[68\%,100\%]$), the highest of any configuration, and 75\% on-target
(CI $[41\%,93\%]$), tied with Fixed-C.

%──────────────────────────────────────────────────────────────────────────────
\section{Conclusions}

\textbf{(1) Fix~C eliminates gradient explosions and achieves the best Fix-A
return ($3221.5$, CV=13.7\%).}
Zero explosions across 8 seeds vs 3/8 for Fixed and Fixed-Warmup. All
$b_\text{final}\in[367,384]$, tightly centred around the target 374, confirming
the cold-start overshoot is eliminated. Six of eight Fixed-C seeds reach the
on-target band ($P(A<0)\in[0.6,0.8]$), the highest rate of any configuration.
However, only seed~4 exceeds $P(A<0)>0.70$ (0.733), and 2/8 seeds have
$\mathbb{E}[A]>0$, so Fix~C does not universally solve the calibration problem.

\textbf{(2) Classical ($n=8$) vs Fix-A pool ($n=24$) is significant ($p=0.019$, $d=0.87$).}
With $n=8$ classical seeds, the classical-vs-Fix-A-pool comparison is significant
($p=0.009$, $d=0.67$), confirming the quantum-classical return gap at the expanded
sample size. Classical mean return $3499.8 \pm 39.5$ vs Fix-A pool $3152 \pm 589$.

\textbf{(3) Fix-A vs default-init on calibration is highly significant ($p=0.0004$).}
Direction-correct rate: Fix-A pool 20/24 (83.3\%) vs default pool 2/11 (18.2\%).
This replicates and strengthens Task 10's finding at substantially larger sample
sizes. The scale initialisation (Fix~A) is the effective component, confirmed
across all three Fix-A variant configurations.

\textbf{(4) All-quantum vs classical is significant on return ($p=0.0007$).}
The aggregate comparison across all 35 quantum seeds vs classical establishes
that quantum configurations collectively underperform on return. The on-target
calibration comparison ($p=0.041$) confirms aggregate calibration deficit.

\textbf{(5) Fixed-C achieves the highest on-target rate (6/8), but Fixed-Warmup
produces the most over-calibrated seeds (seeds 0 and 5 at 0.829 and 0.857).}
The best-calibrated seeds across the Fix-A family are Fixed-Warmup seed~5
($P(A<0)=0.857$) and seed~0 ($0.829$), both over-calibrated. This confirms that
warmup removal is not always beneficial; retaining it alongside Fix~A can
produce the highest calibration values at the cost of layer-boundary explosions.

\textbf{(6) Classical-Deep confirms depth is not the quantum instability source.}
Eight seeds of depth-3 MLP $[8,8,8]$ yield $3307.9\pm413.2$ (CV=12.5\%, driven
by a single anomalous seed). Zero explosions. The DRU circuit's depth-3
effective capacity is comparable to the classical depth-3 baseline; the
instability is in the quantum-specific training dynamics, not the architecture.

\textbf{(7) Quantum MQR achieves best calibration (8/8 direction-correct) but
introduces a new failure mode.}
The multi-qubit readout (8 independent $a_i$ parameters) achieves 100\%
direction-correct and 75\% on-target calibration — the best rates of any quantum
configuration. However, two seeds fail catastrophically (2213, 2375 return) despite
good calibration, with high advantage variance ($\text{adv\_std}\in[23,85]$) likely
disrupting action selection. Mean return $3206.3 \pm 573.4$ is not significantly
different from classical ($p=0.191$) or Fix-C ($p=0.955$).

\textbf{(8) Classical-Deep confirms depth is not the quantum instability source.}
Eight seeds of depth-3 MLP $[8,8,8]$ yield $3307.9\pm413.2$ (CV=12.5\%, driven
by a single anomalous seed). Zero explosions. The DRU circuit's depth-3
effective capacity is comparable to the classical depth-3 baseline; the
instability is in the quantum-specific training dynamics, not the architecture.

\textbf{Future experiments should prioritise:}
\begin{enumerate}
  \item \textbf{Scale the classical baseline to $n=8$} to enable properly
        powered return comparisons with the Fix-A family (currently limited by
        $n=3$ classical).
  \item \textbf{Test on sparse-reward tasks} (e.g.\ \texttt{antmaze-medium-play})
        where reward-weighted behavioural cloning is not competitive. This is the
        critical test of whether Fix-C enables the quantum value network to
        add anything beyond what a constant-$V$ baseline provides.
  \item \textbf{Investigate the affine-head bottleneck}: $V(s)=a\cdot\langle Z_0\rangle+b$
        limits the value surface to a one-dimensional embedding of state. Multi-qubit
        readout or deeper classical post-processing could relieve this constraint.
  \item \textbf{Test Fix-C with longer freeze windows} (e.g.\ 3\,000 steps) to
        determine whether the calibration ceiling can be lowered further.
\end{enumerate}

%──────────────────────────────────────────────────────────────────────────────
\begin{thebibliography}{9}
\bibitem{kostrikov2021}
I.~Kostrikov, A.~Nair, and S.~Levine,
``Offline Reinforcement Learning with Implicit Q-Learning,''
\textit{ICLR}, 2022.

\bibitem{skolik2021}
A.~Skolik, J.~J.~McClean, M.~Mohseni, P.~van der Smagt, and M.~Leib,
``Layerwise learning for quantum neural networks,''
\textit{Quantum Machine Intelligence}, 3(1):5, 2021.
\end{thebibliography}

\end{document}